\documentclass[11pt]{article}

\usepackage[margin=1in]{geometry}
\usepackage{amsmath,amssymb,amsthm,mathtools}
\usepackage{booktabs}
\usepackage{microtype}

\newtheorem{theorem}{Theorem}
\newtheorem{lemma}[theorem]{Lemma}
\newtheorem{proposition}[theorem]{Proposition}
\newtheorem{corollary}[theorem]{Corollary}
\theoremstyle{definition}
\newtheorem{definition}[theorem]{Definition}
\newtheorem{remark}[theorem]{Remark}

\newcommand{\causal}{\mathrm{causal}}
\newcommand{\E}{\mathbb{E}}
\newcommand{\Prob}{\mathbb{P}}
\newcommand{\R}{\mathbb{R}}
\newcommand{\eps}{\varepsilon}
\title{The sharp square-root bill of an exact causal source}
\author{}
\date{3 October 2026}

\begin{document}
\maketitle

\begin{abstract}
Let \(P\) and \(Q\) be finite nondegenerate laws with the same Shannon entropy \(h\) and surprisal standard deviations \(\sigma_P\neq\sigma_Q\).
An exact horizon-\(n\) causal source is a finite seed which, under every deterministic nonanticipative policy, emits the product law of the requested actions.
Its minimal seed entropy admits the second-order expansion
\[
  C_n^{\causal}(P,Q)
  =
  nh+\sqrt{\frac{2}{\pi}}\,|\sigma_P-\sigma_Q|\,\sqrt{n}+o(\sqrt{n}).
\]
The lower bound is not the two-transcript profile bound.
It comes from evaluating the seed atomwise against a grid of threshold policies simultaneously: the seed mass is at most every transcript probability at once, so the seed surprisal dominates the maximum of the centered policy martingales.
Those martingales converge to two-speed Brownian motions, and the lower envelope of their time-\(1\) laws has mean exactly \(\sqrt{2/\pi}\,|\sigma_P-\sigma_Q|\).
The matching upper bound is a one-way quantile recycling of two independent streams, lifted to a history-indexed source by reading the second stream backwards.
The recycled fiber length tracks the running maximum of the surprisal-difference walk, which has the same leading constant.
For the crash pair \(P=(1/2,1/4,1/4)\) and \(H(Q)=3/2\) the coefficient is \(0.0347472540518177\ldots\).
\end{abstract}

\section{Model}

Let \(P\) and \(Q\) be probability laws on finite alphabets, with every atom strictly positive, and write
\[
  h=H(P)=H(Q),\qquad
  \sigma_P^2=\mathrm{Var}(-\log_2 P(X)),\qquad
  \sigma_Q^2=\mathrm{Var}(-\log_2 Q(Y)).
\]
Both variances are finite.
The argument is symmetric in \(P\) and \(Q\); from Section~\ref{sec:threshold} onward we set
\[
  a=\sigma_P\ge b=\sigma_Q
\]
and write \(\Delta=a-b\ge 0\).
The centered one-letter surprisals are uniformly bounded: there is a finite \(L_\star\) such that
\[
  \bigl|-\log_2 P(x)-h\bigr|\le L_\star,\qquad
  \bigl|-\log_2 Q(y)-h\bigr|\le L_\star
\]
for every positive atom.

A deterministic nonanticipative policy \(\pi=(\pi_t)_{t=1}^n\) chooses an action
\[
  A_t=\pi_t(Y_1,\ldots,Y_{t-1})\in\{P,Q\}.
\]
An exact horizon-\(n\) history-indexed causal source is a finite random seed \(S\) together with deterministic output maps such that, for every such policy, the transcript \(Y^\pi\) has law
\begin{equation}
  \label{eq:exact}
  R_\pi(y^n)=\prod_{t=1}^n P_{A_t}(y_t),
\end{equation}
where \(P_P=P\) and \(P_Q=Q\).
Equivalently, \eqref{eq:exact} holds for every deterministic action word, and hence for every adaptive policy.
The causal cost is the minimal seed entropy
\[
  C_n^{\causal}(P,Q)=\inf\bigl\{H(S):S\text{ is an exact horizon-\(n\) seed}\bigr\}.
\]
For a policy \(\pi\) and a transcript \(y^n\) write
\[
  I_\pi(y^n)=-\log_2 R_\pi(y^n),\qquad
  M_{\pi,n}(y^n)=I_\pi(y^n)-nh.
\]
Thus \(M_{\pi,n}(Y^\pi)\) is a centered martingale in the filtration of \(Y^\pi\), with increments bounded by \(L_\star\).

\section{The seed sees every policy at once}
\label{sec:seed}

The two-policy profile bound only recovers half the constant below.
The sharp lower bound uses the stronger atomwise comparison.

\begin{lemma}[Atomwise seed inequality]
\label{lem:atom}
Let \(S\) be an exact seed and let \(\Pi\) be any finite family of deterministic nonanticipative policies.
For every atom \(s\) of mass \(\gamma_s>0\),
\[
  \gamma_s\le R_\pi\bigl(y^\pi(s)\bigr)
  \qquad\text{for every }\pi\in\Pi,
\]
where \(y^\pi(s)\) is the transcript produced by \(\pi\) on the seed atom \(s\).
Consequently
\begin{equation}
  \label{eq:seed-max}
  H(S)\ge nh+\E\Bigl[\max_{\pi\in\Pi} M_{\pi,n}(Y^\pi)\Bigr],
\end{equation}
the expectation being under the seed law.
In particular
\[
  C_n^{\causal}(P,Q)
  \ge
  nh+\sup_{\Pi}\E\Bigl[\max_{\pi\in\Pi} M_{\pi,n}\Bigr].
\]
\end{lemma}

\begin{proof}
The output maps are deterministic, so each policy transcript is a function of \(S\).
Exactness says that the push-forward of the seed law under \(s\mapsto y^\pi(s)\) equals \(R_\pi\).
Therefore, for each transcript \(y\),
\[
  \sum_{s:\,y^\pi(s)=y}\gamma_s=R_\pi(y).
\]
Every summand is nonnegative, so \(\gamma_s\le R_\pi(y^\pi(s))\).
The same atom \(s\) satisfies the inequality for every \(\pi\in\Pi\) simultaneously, hence
\[
  -\log_2\gamma_s
  \ge
  \max_{\pi\in\Pi} I_\pi\bigl(y^\pi(s)\bigr)
  =
  nh+\max_{\pi\in\Pi} M_{\pi,n}\bigl(y^\pi(s)\bigr).
\]
Averaging over \(s\) gives \eqref{eq:seed-max}.
\end{proof}

The inequality is strictly stronger than a bound that uses only marginal transcript laws.
The profile functional of the companion note,
\[
  \mathcal L(Z_{\pi_1},\ldots,Z_{\pi_m})
  =\int_0^\infty \max_i\Prob(Z_{\pi_i}>t)\,dt,
\]
depends on the marginal tails alone.
Lemma~\ref{lem:atom} charges the joint maximum realized by one seed.
Azuma's inequality still shows that each \(M_{\pi,n}\) is \(O(\sqrt{n})\) in \(L^1\), uniformly in the policy, so the gain available here is a constant factor inside the square-root scale, not a change of exponent.

For two policies the profile functional collapses to the familiar Wasserstein form, recorded only for comparison.

\begin{proposition}[Two-policy floor]
\label{prop:two}
For the two constant policies,
\[
  C_n^{\causal}(P,Q)
  \ge
  nh+\tfrac12 W_1(S_{P,n},S_{Q,n}),
\]
where \(S_{P,n}\) and \(S_{Q,n}\) are the product surprisal sums.
Hence
\[
  \liminf_{n\to\infty}
  \frac{C_n^{\causal}(P,Q)-nh}{\sqrt{n}}
  \ge
  \frac{|\sigma_P-\sigma_Q|}{\sqrt{2\pi}}.
\]
\end{proposition}

\begin{proof}
Lemma~\ref{lem:atom} with the two constant policies yields
\(H(S)\ge\E\max\{S_{P,n},S_{Q,n}\}\).
For any coupling of two nonnegative random variables,
\(\E\max(U,V)\ge\bigl(\E U+\E V+W_1(\mathcal L(U),\mathcal L(V))\bigr)/2\),
and both transcript entropies equal \(nh\).
Finite support gives \(W_1\) convergence of the normalized sums to the centered Gaussians of variances \(\sigma_P^2\) and \(\sigma_Q^2\).
The monotone coupling of those Gaussians has cost \(|\sigma_P-\sigma_Q|\,\E|G|=\sqrt{2/\pi}\,|\sigma_P-\sigma_Q|\), and the factor \(1/2\) produces the stated liminf.
\end{proof}

The constant in Proposition~\ref{prop:two} is exactly half the constant in Theorem~\ref{thm:main}.
The rest of the lower bound is the recovery of that factor two.

\section{Threshold policies}
\label{sec:threshold}

Fix \(z\in\R\).
The threshold policy \(\pi_{z,n}\) requests the high-variance action while the centered surprisal accumulated so far lies below the barrier, and the low-variance action afterwards:
\[
  A_t=
  \begin{cases}
    P & \text{if }M_{t-1}<z\sqrt{n},\\
    Q & \text{if }M_{t-1}\ge z\sqrt{n}.
  \end{cases}
\]
Write \(M^{(z)}_n\) for its terminal centered surprisal and \(X^{(z)}\) for the diffusion
\begin{equation}
  \label{eq:sde}
  dX_t=\sigma_z(X_t)\,dB_t,
  \qquad
  X_0=0,
  \qquad
  \sigma_z(x)=
  \begin{cases}
    a & \text{if }x<z,\\
    b & \text{if }x\ge z.
  \end{cases}
\end{equation}
The diffusion coefficient is Borel, bounded, and uniformly elliptic.
The martingale problem for \eqref{eq:sde} is well posed, and the time-\(1\) law is absolutely continuous
\cite{Lejay2006,EtoreMartinez}.
Call this time-\(1\) law \(F_z\).

\begin{lemma}[Diffusion limit of one threshold]
\label{lem:fclt}
For each fixed \(z\),
\[
  \frac{M^{(z)}_n}{\sqrt{n}}\ \Rightarrow\ X^{(z)}_1.
\]
The convergence holds jointly for any finite set of thresholds.
Moreover the family \(\{|M^{(z)}_n|/\sqrt{n}\}\) is uniformly integrable, uniformly in \(z\) and \(n\).
\end{lemma}

\begin{proof}
The centered surprisal is a martingale with bounded increments.
Conditional on the past, the next increment has mean zero and variance \(a^2\) or \(b^2\) according to the side of the barrier \(z\sqrt{n}\) on which the current position lies; the conditional variance equals the exact one-letter varentropy of the requested action.
Standard functional convergence for triangular martingale arrays with bounded increments and state-dependent conditional variance applies: the normalized path converges weakly in Skorokhod space to the unique solution of the martingale problem \eqref{eq:sde}.
See Ethier--Kurtz, Chapter~7, for the martingale-problem criterion, and Lejay for uniqueness at a single discontinuous interface.
Joint convergence on a finite grid follows by applying the same criterion to the vector of policies driven by a common seed, or, for the marginal envelope argument below, is unnecessary: only marginal convergence is required.

Uniform integrability is Azuma's inequality.
The increments are bounded by \(L_\star\), so for every policy and every \(s>0\),
\[
  \Prob\bigl(|M_{\pi,n}|>s\bigr)
  \le
  2\exp\Bigl(-\frac{s^2}{2n L_\star^2}\Bigr).
\]
Hence \(\E[|M_{\pi,n}|/\sqrt{n}]^{1+\delta}\) is bounded for every \(\delta\in(0,1)\), uniformly in the policy.
\end{proof}

\section{The envelope}
\label{sec:envelope}

Let \(\Phi\) be the standard normal distribution function.
Define
\begin{equation}
  \label{eq:Fstar}
  F_*(x)=
  \begin{cases}
    \dfrac{2b}{a+b}\,\Phi(x/b) & \text{if }x\le 0,\\[1.2ex]
    1-\dfrac{2a}{a+b}\bigl(1-\Phi(x/a)\bigr) & \text{if }x>0.
  \end{cases}
\end{equation}
This is a distribution function.
Its mean is the constant we need, and the computation is elementary once \eqref{eq:Fstar} is granted.

\begin{lemma}[Mean of the envelope]
\label{lem:mean}
If \(R_*\) has law \(F_*\), then
\[
  \E[R_*]=\sqrt{\frac{2}{\pi}}\,(a-b).
\]
\end{lemma}

\begin{proof}
The Gaussian tail integrals
\[
  \int_0^\infty\bigl(1-\Phi(x/c)\bigr)\,dx=\frac{c}{\sqrt{2\pi}},
  \qquad c>0,
\]
give
\begin{align*}
  \E[R_*]
  &=
  \int_0^\infty\bigl(1-F_*(x)\bigr)\,dx
  -\int_{-\infty}^0 F_*(x)\,dx\\
  &=
  \frac{2a}{a+b}\cdot\frac{a}{\sqrt{2\pi}}
  -\frac{2b}{a+b}\cdot\frac{b}{\sqrt{2\pi}}
  =
  \frac{2}{\sqrt{2\pi}}\,(a-b)
  =
  \sqrt{\frac{2}{\pi}}\,(a-b).
\end{align*}
The identity \(a^2-b^2=(a-b)(a+b)\) is the only cancellation.
\end{proof}

The distribution \(F_*\) is not the law of a single two-speed path.
It is the lower envelope of the family \(\{F_z\}\), attained by placing the interface at the level being tested.

\begin{lemma}[Interface formula]
\label{lem:interface}
For every \(x\in\R\),
\[
  F_*(x)=\Prob\bigl(X^{(x)}_1\le x\bigr)
  =\inf_{z\in\R}\Prob\bigl(X^{(z)}_1\le x\bigr).
\]
In particular \(F_*(x)=\inf_z F_z(x)\).
\end{lemma}

\begin{proof}
Write \(u(t,y)=\Prob(X_1\le x\mid X_t=y)\) for a controlled diffusion whose instantaneous volatility \(\sigma_t\) is progressively measurable and takes values in \(\{a,b\}\).
Then \(u(1,y)=\mathbf 1_{y\le x}\), and on \([0,1)\) the function \(u\) satisfies the HJB equation
\begin{equation}
  \label{eq:hjb}
  \partial_t u+\tfrac12\sup_{\sigma\in\{a,b\}}\sigma^2\,\partial_{yy}u=0
\end{equation}
in the viscosity sense, because the controller who wants to minimize \(\Prob(X_1\le x)\) chooses the volatility that the forward equation charges with a minus sign; equivalently, the probability of the bad half-line is the expectation of the indicator under the controlled law, and the generator contribution of a pure diffusion is \(\tfrac12\sigma^2\partial_{yy}\).
The candidate feedback
\[
  \sigma^*(y)=a\cdot\mathbf 1_{y<x}+b\cdot\mathbf 1_{y\ge x}
\]
is exactly \(\pi_{x}\) in the diffusion limit.
Its occupation formula is classical for oscillating Brownian motion after a translation that moves the interface to the origin
\cite{KeilsonWellner,Lejay2006}.
Translated back, the sub-probability density of \(X^{(x)}_1\) on \((-\infty,x]\) is the left-hand Gaussian density of variance \(b^2\), weighted by the transmission factor \(2b/(a+b)\) when the starting point lies on the low-volatility side of the interface, and the right-hand tail is the Gaussian density of variance \(a^2\), weighted by \(2a/(a+b)\).
Those are the two lines of \eqref{eq:Fstar}, so the candidate achieves \(F_*\).

To see that no other interface does better, observe that \(F_*\) is a viscosity solution of \eqref{eq:hjb} with the same terminal condition.
On \(\{y<x\}\) the candidate value function is a mixture of Gaussians whose second derivative has the sign that selects \(\sigma=a\); on \(\{y>x\}\) it selects \(\sigma=b\).
Uniqueness of bounded viscosity solutions of \eqref{eq:hjb} gives optimality of the threshold \(z=x\).
Thus the pointwise infimum over deterministic thresholds, and even over all progressively measurable volatility choices, equals \(F_*(x)\).
\end{proof}

At the single point \(x=0\) the formula reduces to the transmission weight
\[
  \Prob\bigl(X^{(0)}_1\le 0\bigr)=\frac{b}{a+b},
\]
which is the standard left-mass of oscillating Brownian motion started on the interface.
For \(x\neq 0\) the starting point is not on the interface, and the tail scale is the volatility of the side that contains the origin: scale \(b\) for \(x<0\), scale \(a\) for \(x>0\).

\begin{lemma}[Envelope lower bound]
\label{lem:envelope-exp}
Let \(\Pi_n\) be a finite \(n^{-1/4}\)-grid of thresholds in \([-\log n,\log n]\).
Then
\[
  \liminf_{n\to\infty}
  \E\Bigl[\max_{\pi\in\Pi_n}\frac{M_{\pi,n}}{\sqrt{n}}\Bigr]
  \ge
  \sqrt{\frac{2}{\pi}}\,(a-b).
\]
\end{lemma}

\begin{proof}
Write \(M_n=\max_{\pi\in\Pi_n} M_{\pi,n}/\sqrt{n}\).
For every \(x\),
\[
  \Prob(M_n\le x)
  \le
  \min_{\pi\in\Pi_n}\Prob(M_{\pi,n}/\sqrt{n}\le x).
\]
Pass to the limit along the grid.
Lemma~\ref{lem:fclt} gives marginal convergence at every fixed threshold.
The map \(z\mapsto F_z(x)\) is continuous in \(z\) at continuity points of \(F_z\), because the diffusion coefficient depends continuously on the interface in the bounded-Lipschitz metric on time-\(1\) laws: moving the interface by \(\eta\) changes the volatility on an interval of length \(\eta\), and the resulting total-variation distance is \(O(\eta^{1/2})\) by the Burkholder--Davis--Gundy inequality.
A grid of spacing \(n^{-1/4}\) is therefore dense enough that
\[
  \limsup_{n\to\infty}\Prob(M_n\le x)
  \le
  \inf_z F_z(x)=F_*(x)
\]
at every continuity point of \(F_*\).
Layer-cake representation and Fatou's lemma on the positive part give
\[
  \liminf_{n\to\infty}\E[(M_n)_+]
  \ge
  \int_0^\infty\bigl(1-F_*(t)\bigr)\,dt.
\]
For the negative part the inequality runs the other way and produces a matching limsup,
\[
  \limsup_{n\to\infty}\E[(M_n)_-]
  \le
  \int_{-\infty}^0 F_*(t)\,dt,
\]
after truncating at \(\log n\) and using the Azuma tail to show that the truncation costs \(o(1)\).
Thresholds outside \([-\log n,\log n]\) do not contribute to the envelope at the \(\sqrt{n}\) scale: a barrier beyond \(\log n\) standard deviations is not hit with probability \(1-o(n^{-K})\) for every \(K\), again by Azuma, so the grid restriction is free.
Combining the two sides and applying Lemma~\ref{lem:mean} yields the claim.
Uniform integrability of the maximum over the grid follows from the same Azuma tail and a union bound: the grid has cardinality \(O(n^{1/4}\log n)\), which does not damage the Gaussian tail.
\end{proof}

\section{Matching upper bound}
\label{sec:upper}

The construction has two layers.
First produce a one-way joint of two length-\(n\) strings whose conditional entropy meets the constant.
Then lift that joint to an exact history-indexed seed by reading the second string backwards.

\subsection{One-way recycling}

Let \(L=\lfloor(\log n)^2\rfloor\) and \(m=\lfloor n/L\rfloor\), and write \(n=mL+r\) with \(0\le r<L\).
The remainder \(r\) contributes \(O(L)=o(\sqrt{n})\) bits and is ignored after the first sentence of the entropy budget.

Draw \(W^n\sim Q^n\).
Generate \(X^n\) from \(W^n\) and an auxiliary uniform seed by blocks, causally.
On block \(j=1,\ldots,m\), having seen \(W\) on that block, split the current seed fiber in the proportions of \(P^{\otimes L}\) and retain the atom consistent with a \(P\)-block that is quantile-coupled to the observed \(Q\)-block.
Concretely, let \(U_0\sim\mathrm{Unif}[0,1]\) be independent of \(W^n\), and maintain the dyadic fiber \((\lambda_j^-,\lambda_j^+)\) containing \(U_0\).
The \(Q\)-block \(W^{(j)}\) has surprisal \(B_j=-\log_2 Q^{\otimes L}(W^{(j)})\).
The comonotone (quantile) preimage of this block under the product law \(P^{\otimes L}\) is an interval of \(P^{\otimes L}\)-mass
\[
  \alpha_j=\Theta\bigl(2^{-(B_j-A_j)}\bigr),
\]
where \(A_j\) is the \(P\)-surprisal of the selected atom; the implicit factor is the usual arithmetic-coding redundancy, at most \(2\) per block after the Kraft tightening of \(O(1)\) bits per block recorded below.
Intersecting the current fiber with that preimage multiplies its length by at most \(\alpha_j\).
The resulting string \(X^n\) is exactly \(P^n\)-distributed, because each block is drawn from the \(P^{\otimes L}\) splitting of a uniform fiber, and \(X_t\) is a function of \(U_0\) and \(W_1,\ldots,W_t\) only.

\begin{lemma}[Fiber length tracks the running maximum]
\label{lem:fiber}
Let \(A_j,B_j\) be the block surprisals under the comonotone coupling, and set
\[
  Z_j=(A_j-Lh)-(B_j-Lh),\qquad
  S_k=\sum_{j=1}^k Z_j,\qquad
  S_k^+=\max_{0\le i\le k} S_i
\]
with \(S_0=0\).
The terminal fiber length \(\lambda\) satisfies
\begin{equation}
  \label{eq:fiber}
  -\log_2\lambda
  \le
  Lh\cdot 0+S_m^++G_m
  =
  S_m^++G_m,
\end{equation}
where the overhead \(G_m\) obeys \(G_m\le C\log m\) almost surely for a constant \(C\) depending only on the alphabets.
Consequently
\[
  H(X^n\mid W^n)
  \le
  \E[S_m^+]+O(\log n).
\]
\end{lemma}

\begin{proof}
The construction never needs more uniform bits than the heaviest fiber restriction encountered along the path.
A negative increment \(Z_j<0\) shortens the \(P\)-preimage relative to the \(Q\)-surprisal already paid by \(W\), and the unused portion of the fiber can be returned to the pool: the nesting is a restriction followed by a release, not a permanent multiplication by every \(\alpha_j\).
The outstanding restriction at step \(k\) is therefore bounded by the largest partial sum of the increments \(Z_j\), up to the integer-bit loss of an arithmetic coder.
Each block contributes at most \(O(1)\) bits of coder redundancy (a fixed Kraft gap coming from the finite alphabet), and the record times of \(S_k\) contribute an additional \(O(\log m)\) bits to address the active interval endpoints.
This is the inequality \eqref{eq:fiber}.
The conditional entropy is at most the expected codelength of a prefix-free encoding of the fiber, which is \(-\E\log_2\lambda+O(1)\).
\end{proof}

The overhead \(O(\log n)\) is \(o(\sqrt{n})\).
The walk itself produces the constant.

\begin{lemma}[Walk maximum]
\label{lem:walk}
Under the comonotone coupling of the two block-surprisal laws,
\[
  \E[Z_1]=0,
  \qquad
  \frac{\E[Z_1^2]}{L}\to(a-b)^2,
\]
and
\[
  \frac{\E[S_m^+]}{\sqrt{n}}
  \to
  \sqrt{\frac{2}{\pi}}\,(a-b).
\]
\end{lemma}

\begin{proof}
Each block is an i.i.d.\ copy, so \((Z_j)_{j\le m}\) is an i.i.d.\ centered triangular array.
Finite alphabets give a fourth-moment bound \(\E Z_1^4=O(L^2)\).
The normalized block sums obey a CLT, and the comonotone coupling is the \(W_2\)-optimal coupling of the two block laws.
Therefore
\[
  \frac{Z_1}{\sqrt{L}}
  \Rightarrow
  (a-b)G,
  \qquad G\sim N(0,1),
\]
which is the variance claim; the fourth-moment bound upgrades convergence of second moments.
Donsker's theorem lifts the array to a Brownian path on \([0,1]\) after diffusive scaling by \(\sqrt{mL}=\sqrt{n-r}\).
The running-maximum functional is continuous in the Skorokhod topology on the set of paths whose maximum is not attained at a flat stretch, a set of full Wiener measure.
Hence \(S_m^+/\sqrt{n}\) converges in law to \((a-b)\max_{t\le 1} B_t\).
Uniform integrability follows from Doob's inequality and the fourth-moment bound.
The reflected Brownian maximum has expectation
\[
  \E\Bigl[\max_{t\le 1} B_t\Bigr]=\sqrt{\frac{2}{\pi}},
\]
by the reflection identity \(\Prob(\max_{t\le 1} B_t\ge u)=2\Prob(B_1\ge u)\) for \(u\ge 0\).
Multiplying by \(a-b\) gives the limit.
\end{proof}

\subsection{Anti-diagonal lift}

It remains to turn the one-way joint into an exact causal seed without increasing the entropy by more than \(o(\sqrt{n})\).

\begin{lemma}[Anti-diagonal lift]
\label{lem:lift}
Suppose \(X^n\sim P^n\), \(W^n\sim Q^n\), and
\begin{equation}
  \label{eq:oneway}
  \mathcal L(X_{1:k}\mid W_{1:n})
  =
  \mathcal L(X_{1:k}\mid W_{1:k})
  \qquad\text{for every }k\le n.
\end{equation}
Read \(W\) backwards: on the \(r\)-th request for \(P\), emit \(X_r\); on the \(s\)-th request for \(Q\), emit \(W_{n-s+1}\).
Every action word, and every adaptive policy, then exposes independent symbols with the requested product law.
The pair \((X^n,W^n)\) is an exact horizon-\(n\) seed, and
\[
  C_n^{\causal}(P,Q)
  \le
  H(X^n,W^n)
  =
  nh+H(X^n\mid W^n).
\]
\end{lemma}

\begin{proof}
An action word with \(r\) letters \(P\) and \(s\) letters \(Q\), \(r+s=n\), exposes \(X_1,\ldots,X_r\) and \(W_n,W_{n-1},\ldots,W_{n-s+1}\).
The reversed block is exactly \(W_{r+1},\ldots,W_n\), hence disjoint from \(\{1,\ldots,r\}\).
Condition \eqref{eq:oneway} says that \(X_{1:r}\) is a function of the seed fiber and \(W_{1:r}\) only, so it is independent of \(W_{r+1:n}\).
The marginals are the requested products by construction.
An adaptive policy reveals only a prefix of some action word at each time, and the same index disjointness applies pathwise: after \(r\) requests for \(P\) and \(s\) requests for \(Q\), the reversed reads still lie strictly above index \(r\).
Thus every policy transcript has law \(R_\pi\).
\end{proof}

No padding is required, and a padding segment would be illegal for the entropy budget: an unused middle copy of \(W\) adds a second \(nh\) and destroys the first-order term.
The reversed block of a mixed word already occupies the complementary index set \(\{r+1,\ldots,n\}\), so the same string serves both directions.
In the recycling construction \(X_t\) is a function of the fiber of \(U_0\) and of \(W_{1:t}\), which is \eqref{eq:oneway}.
The seed stores \(W^n\) and a prefix-free encoding of the occupied fiber.
Lemma~\ref{lem:fiber} bounds that fiber, and therefore
\[
  C_n^{\causal}(P,Q)
  \le
  nh+\E[S_m^+]+O(\log n).
\]

\begin{lemma}[Upper bound]
\label{lem:upper}
\[
  \limsup_{n\to\infty}
  \frac{C_n^{\causal}(P,Q)-nh}{\sqrt{n}}
  \le
  \sqrt{\frac{2}{\pi}}\,|\sigma_P-\sigma_Q|.
\]
\end{lemma}

\begin{proof}
Combine the implementation just given with Lemmas~\ref{lem:fiber} and~\ref{lem:walk}.
The block length \(L=(\log n)^2\) tends to infinity, so the CLT variance identification applies, and \(L=o(n^{1/2})\) keeps the edge block inside the \(o(\sqrt{n})\) remainder.
\end{proof}

\section{The theorem}

\begin{theorem}
\label{thm:main}
Let \(P,Q\) be finite strictly positive laws with \(H(P)=H(Q)=h\) and surprisal standard deviations \(\sigma_P,\sigma_Q\).
Then
\begin{equation}
  \label{eq:main}
  C_n^{\causal}(P,Q)
  =
  nh+\sqrt{\frac{2}{\pi}}\,|\sigma_P-\sigma_Q|\,\sqrt{n}+o(\sqrt{n}).
\end{equation}
If \(\sigma_P=\sigma_Q\), the second term vanishes and the construction of Section~\ref{sec:upper} gives \(C_n^{\causal}=nh+O(\log n)\).
\end{theorem}

\begin{proof}
The case \(\Delta=0\) is the upper bound with a degenerate walk, whose maximum is \(O(\log n)\) by the fourth-moment bound, together with the trivial \(C_n^{\causal}\ge nh\).
For \(\Delta>0\), Lemma~\ref{lem:atom} and Lemma~\ref{lem:envelope-exp} give the liminf, and Lemma~\ref{lem:upper} gives the limsup.
\end{proof}

The profile floor of Proposition~\ref{prop:two} is the same statement with \(\sqrt{2/\pi}\) replaced by \(1/\sqrt{2\pi}\).
The factor two is the difference between a terminal Gaussian comparison and a running maximum: \(\E[(B_1)_+]=\sqrt{1/(2\pi)}\) and \(\E[\max_{t\le 1} B_t]=\sqrt{2/\pi}\).

\section{Crash pair}

Take \(P=(1/2,1/4,1/4)\) and \(Q=(q,q,1-2q)\), with \(q\in(1/3,1/2)\) the unique solution of \(H(Q)=3/2\).
Then
\begin{align*}
  q&=0.4100324281819856\ldots,\\
  h&=3/2,\\
  \sigma_P^2&=1/4,\\
  \sigma_Q^2&=0.2083473102390734\ldots,\\
  \sigma_P&=1/2,\\
  \sigma_Q&=0.4564507752639635\ldots.
\end{align*}
Theorem~\ref{thm:main} specializes to
\begin{equation}
  \label{eq:crash}
  C_n^{\causal}(P,Q)
  =
  \tfrac32 n
  +0.0347472540518177\ldots\,\sqrt{n}
  +o(\sqrt{n}),
\end{equation}
since
\[
  \sqrt{\frac{2}{\pi}}\,|\sigma_P-\sigma_Q|
  =0.0347472540518177\ldots.
\]
The two-policy floor is half of this, \(0.0173736270259089\ldots\).
The independent antidiagonal block construction remains a valid, coarser upper bound of order \(n^{2/3}\), with leading coefficient \(0.145119355153\ldots\); Theorem~\ref{thm:main} closes that gap.

\section{What the argument uses, and what it does not}

The atomwise comparison, Lemma~\ref{lem:atom}, is elementary and does not pass through the profile functional \(\mathcal L\).
The square-root barrier for \(\mathcal L\) proved in the companion note is therefore not a barrier for \eqref{eq:seed-max}.
That barrier says a universal \(n^{2/3}\) lower bound cannot be extracted from marginal transcript tails.
It does not forbid a sharp constant on the \(\sqrt{n}\) scale, which is all that Theorem~\ref{thm:main} claims.

The analytic inputs are standard and localized.
Well-posedness of \eqref{eq:sde} is the martingale problem for a uniformly elliptic piecewise constant diffusion coefficient \cite{Lejay2006}.
The passage from marginal convergence to the expectation of the maximum is the envelope inequality \(\Prob(\max_z R_z\le x)\le\inf_z\Prob(R_z\le x)\), which holds for every coupling, plus Azuma tails for uniform integrability.
No identification of the joint law of \((X^{(z)})_z\) is required on the lower side.
On the upper side the blocks are independent by construction, so there is no dependence-removal step: \(L=(\log n)^2\) is used only to reach the Gaussian variance \((a-b)^2 L\).
Prefix-free exactization of \(U_0\) costs \(O(\log n)\) bits.
The one-way condition \eqref{eq:oneway} is preserved because each output symbol is a function of the seed fiber and the past of \(W\) only, and the reversed \(Q\)-reads lie off that past.

\begin{thebibliography}{9}
\bibitem{EtoreMartinez}
P.~\'Etor\'e and M.~Martinez,
\emph{On the existence of a time-inhomogeneous skew Brownian motion and some related laws},
Electron.\ J.\ Probab.\ \textbf{17} (2012), no.~19.
\bibitem{KeilsonWellner}
J.~Keilson and J.~A.~Wellner,
\emph{Oscillating Brownian motion},
J.\ Appl.\ Probab.\ \textbf{15} (1978), 300--310.
\bibitem{Lejay2006}
A.~Lejay,
\emph{On the constructions of the skew Brownian motion},
Probab.\ Surv.\ \textbf{3} (2006), 413--466.
\end{thebibliography}

\end{document}
